\chapter{Componentes conexas, caminos y arcos}\label{chap:18}

\section{De la conexidad global a sus piezas maximales}
Un espacio disconexo no siempre se descompone en abiertos conexos. Las componentes proporcionan sus piezas conexas maximales, que son siempre cerradas pero pueden dejar de ser abiertas. Por otra parte, unir dos puntos mediante una curva es una exigencia más concreta que la conexidad. Conviene distinguir un camino continuo de un arco inyectivo, pues la concatenación conserva continuidad pero puede introducir repeticiones. Este capítulo desarrolla la lectura 18 de \cite{valencia}, usando las propiedades de conexidad del capítulo \ref{chap:17}.

\section{Componentes y desconexión total}
\begin{proposicion}\label{ch18:componentes}
En un espacio métrico $X$, definamos $x\sim y$ si existe un subconjunto conexo que contiene a ambos. Esta es una relación de equivalencia. Sus clases, llamadas \emph{componentes conexas}, son precisamente los subconjuntos conexos no vacíos maximales por inclusión. Todas ellas son cerradas.\index{componente conexa}
\end{proposicion}
\begin{proof}
La reflexividad se obtiene del conexo $\{x\}$; la simetría es inmediata porque contener a $x,y$ no depende del orden. Si $x,y\in A$ y $y,z\in B$ con $A,B$ conexos, su intersección contiene $y$, así que $A\cup B$ es conexo y contiene $x,z$. Esto prueba transitividad.

La clase $C_x$ de $x$ es la unión de todos los conjuntos conexos que contienen a $x$. Todos comparten ese punto, luego la unión es conexa. Si un conexo $D$ intersecta $C_x$, la unión $D\cup C_x$ es conexa y contiene $x$; por la definición de la clase, toda la unión está en $C_x$, así que $D\subset C_x$. En particular, no puede ampliarse $C_x$ a un conexo mayor. Recíprocamente, todo conexo no vacío $D$ está contenido en la clase de cualquiera de sus puntos; si es maximal, coincide con esa clase.

Finalmente $\cl{C_x}$ es conexo por el ejercicio 3 del capítulo \ref{chap:17}. Como contiene a $C_x$, su maximalidad obliga a $\cl{C_x}=C_x$. Esto demuestra la cerradura.
\end{proof}

\begin{definicion}
Un espacio es \emph{totalmente disconexo}\index{desconexion total@desconexión total} si sus componentes son singletons; equivalentemente, todo subconjunto conexo no vacío consta de un punto.
\end{definicion}
Un espacio no vacío es conexo si y solo si tiene una componente. En un espacio discreto las componentes son puntos, pero el recíproco falla: $\Q$ es totalmente disconexo sin ser discreto. En efecto, si un subconjunto de $\Q$ contiene racionales $a<b$, elegimos un irracional $c\in(a,b)$; sus intersecciones con $(-\infty,c)$ y $(c,\infty)$ lo separan. Además ningún racional está aislado, porque cualquier intervalo a su alrededor contiene otros racionales.

\begin{proposicion}\label{ch18:abiertosEuclideos}
Todo abierto de $\R^n$ es unión de a lo sumo numerables abiertos conexos disjuntos, sus componentes.
\end{proposicion}
\begin{proof}
Sea $A$ abierto y $C$ una componente de $A$. Para $x\in C$, existe una bola $B(x,r)\subset A$. La bola es convexa, luego conexa: cada punto puede unirse con el centro por un segmento y la unión de esos segmentos es conexa. Como la bola intersecta $C$, su unión con $C$ es conexa y la maximalidad implica $B(x,r)\subset C$. Por tanto cada componente es abierta en $\R^n$. Cada una contiene algún punto de $\Q^n$, pues este conjunto es denso. Fijando una enumeración de $\Q^n$, asignamos a cada componente el primer racional de la enumeración que contenga; componentes distintas reciben racionales distintos, así que hay a lo sumo numerables. Son cerradas en $A$ por la proposición anterior, aunque no necesariamente cerradas en $\R^n$: el intervalo $(0,1)$ es una componente de sí mismo.
\end{proof}

\section{Caminos, arcos y el problema de las repeticiones}
\begin{definicion}
Un \emph{camino}\index{camino} en $X$ es una aplicación continua $\gamma:[0,1]\to X$. Un \emph{arco}\index{arco} es un camino inyectivo. Sus extremos son $\gamma(0),\gamma(1)$. Un espacio es \emph{conexo por caminos} si cada pareja de puntos es extremo de un camino. Es \emph{arcoconexo}\index{arcoconexidad} si cada pareja de puntos \emph{distintos} es extremo de un arco. Por convención, un punto está conectado consigo mismo; no se exige un mapa inyectivo con extremos iguales, pues eso sería imposible. El vacío y el singleton cumplen ambas condiciones con estas convenciones.
\end{definicion}
Un arco es un homeomorfismo de $[0,1]$ sobre su imagen, porque es una aplicación continua inyectiva desde un compacto a un espacio métrico. Si dos caminos tienen un extremo común, su concatenación, recorriendo cada uno en la mitad del tiempo, es continua por el lema de pegado sobre los intervalos cerrados $[0,1/2]$ y $[1/2,1]$. Pero incluso si ambos caminos son inyectivos, la concatenación puede visitar un punto dos veces. Para conservar un arco, necesitamos el siguiente recorte.

\begin{lema}[Recorte de dos arcos]\label{ch18:recorte}
Si hay un arco de $x$ a $y$ y otro de $y$ a $z$, entonces, cuando $x\ne z$, hay un arco de $x$ a $z$ contenido en la unión de sus imágenes.
\end{lema}
\begin{proof}
Sean $p,q:[0,1]\to X$ los arcos con $p(0)=x$, $p(1)=q(0)=y$ y $q(1)=z$. Si $z\in p([0,1])$, la inyectividad da un único $s\in(0,1]$ con $p(s)=z$, y $t\mapsto p(st)$ es el arco buscado.

Si $z\notin p([0,1])$, el conjunto
\[
E=\{t\in[0,1]:q(t)\in p([0,1])\}
\]
es cerrado, porque la imagen compacta de $p$ es cerrada. Es no vacío, pues contiene $0$, y no contiene $1$. Tiene máximo $t_0<1$. Existe un único $s_0\in[0,1]$ con $p(s_0)=q(t_0)$. Por maximalidad de $t_0$, la imagen de $q$ para $t>t_0$ no toca la imagen de $p$.

Si $s_0=0$, basta el arco $t\mapsto q(t_0+(1-t_0)t)$. Si $s_0>0$, usamos
\[
r(t)=\begin{cases}
p(2s_0t),&0\le t\le1/2,\\
q(t_0+(2t-1)(1-t_0)),&1/2\le t\le1.
\end{cases}
\]
Las dos expresiones coinciden en $t=1/2$, así que $r$ es continua. Cada tramo es inyectivo y sus imágenes solo se intersectan en el punto de unión, por la elección de $t_0$. Por ello $r$ es inyectiva y tiene extremos $x,z$. Este último control de la intersección es lo que una concatenación sin recorte no garantiza.
\end{proof}

\begin{proposicion}\label{ch18:arcoimplica}
Todo espacio conexo por caminos, y en particular todo espacio arcoconexo, es conexo.
\end{proposicion}
\begin{proof}
Si existiera una separación $U,V$, elegiríamos $x\in U$ y $y\in V$. Un camino $\gamma$ entre ambos tendría imagen conexa por ser imagen continua del intervalo $[0,1]$. Sin embargo, $\gamma([0,1])\cap U$ y $\gamma([0,1])\cap V$ serían una separación no trivial de esa imagen. La contradicción excluye la separación de $X$.
\end{proof}

\begin{ejemplo}[Una curva límite conexa sin conexión por caminos]\label{ch18:seno}
Consideremos
\[
S=\{(x,\sin(1/x)):x\ne0\}\cup(\{0\}\times[-1,1]).
\]
La gráfica $\Gamma_+$ para $x>0$ es conexa por ser imagen continua del intervalo $(0,\infty)$; la gráfica $\Gamma_-$ para $x<0$ también lo es. La cerradura de cada una incluye el segmento vertical $L=\{0\}\times[-1,1]$. Para verlo, dado $y\in[-1,1]$, elegimos $\alpha$ con $\sin\alpha=y$ y puntos $x_n=1/(2\pi n+\alpha)>0$ para $n$ grande. Entonces $(x_n,\sin(1/x_n))\to(0,y)$. Para la rama negativa usamos un ángulo $\beta$ con $-\sin\beta=y$ y $x_n=-1/(2\pi n+\beta)$. No hay otros puntos nuevos en las cerraduras: si la primera coordenada límite no es cero, la continuidad de $\sin(1/x)$ determina la segunda; si es cero, la segunda queda en $[-1,1]$. Por tanto $\cl{\Gamma_+}=\Gamma_+\cup L$ y $\cl{\Gamma_-}=\Gamma_-\cup L$. Son conexos y se intersectan en $L$, de modo que su unión $S$ es conexa.

No existe siquiera un camino en $S$ desde $(1,\sin1)$ hasta $(-1,-\sin1)$. Supongamos que $f(t)=(u(t),v(t))$ fuera uno. Por el valor intermedio, $u$ se anula; sea
\[
t_0=\min\{t\in[0,1]:u(t)=0\}.
\]
El mínimo existe porque el conjunto de ceros es cerrado no vacío en un compacto. Además $t_0>0$ y $u(t)>0$ para $0\le t<t_0$: no puede cambiar de signo sin anularse. Fijemos $0<\eta<t_0$. Tenemos $u(t_0-\eta)>0$. Existen números $a,b$ entre $0$ y $u(t_0-\eta)$ con $\sin(1/a)=1$ y $\sin(1/b)=-1$, tomando denominadores suficientemente grandes de la forma $\pi/2+2\pi n$ y $3\pi/2+2\pi n$. El valor intermedio proporciona $s,r\in(t_0-\eta,t_0)$ con $u(s)=a$, $u(r)=b$. Puesto que sus primeras coordenadas son positivas y $f$ toma valores en $S$, se sigue $v(s)=1$ y $v(r)=-1$.

Esto ocurre para todo $\eta>0$ suficientemente pequeño, así que $v$ no tiene límite por la izquierda en $t_0$, contradiciendo continuidad. El argumento no presupone que los tiempos de una rama y otra formen intervalos, ni identifica límites de composiciones sin justificar qué valores recorre la primera coordenada.
\end{ejemplo}

\section{Abiertos conexos de espacios normados}
\begin{definicion}
En un espacio vectorial real o complejo, el \emph{segmento}\index{segmento} entre $x,y$ es
\[
[x,y]=\{(1-t)x+ty:0\le t\le1\}.
\]
Un conjunto es convexo si contiene el segmento entre cualesquiera dos de sus puntos. Si $x\ne y$, la parametrización del segmento es inyectiva, porque $(t-s)(y-x)=0$ implica $t=s$.
\end{definicion}

\begin{teorema}\label{ch18:abiertoconexo}
Todo abierto conexo de un espacio normado es arcoconexo. En particular, la afirmación vale en un espacio de Banach, sin que se necesite su completitud.
\end{teorema}
\begin{proof}
El vacío y los singletons no presentan dificultad. Fijemos $x\in A$ y sea $P$ el conjunto de puntos que son iguales a $x$ o pueden unirse a $x$ mediante un arco contenido en $A$. Mostraremos que $P$ y $A\setminus P$ son abiertos relativos en $A$.

Si $y\in P$, elegimos $r>0$ con $B(y,r)\subset A$. Para $z\in B(y,r)$, el segmento $[y,z]$ está en esa bola, ya que
\[
\norm{(1-t)y+tz-y}=t\norm{z-y}<r.
\]
Si $y=x$, este segmento proporciona directamente la conexión a $z$; si $y\ne x$, combinamos el arco de $x$ a $y$ con el segmento de $y$ a $z$ y aplicamos el lema de recorte \ref{ch18:recorte}. Los casos de puntos iguales se resuelven por la convención de conexión consigo mismo. Así toda la bola está en $P$ y $P$ es abierto.

Si $y\in A\setminus P$, elegimos otra vez $r>0$ con $B(y,r)\subset A$. Si algún $z$ de esa bola perteneciera a $P$, podríamos unir $x$ a $z$ y después unir $z$ a $y$ por un segmento dentro de la bola; el lema de recorte produciría un arco de $x$ a $y$, contradicción. Por tanto la bola está en $A\setminus P$, y este conjunto también es abierto. Como $A$ es conexo y $x\in P$, debe ser $P=A$. Al haber elegido arbitrariamente el punto inicial $x$, cualesquiera dos puntos distintos de $A$ se unen por un arco.
\end{proof}

\section{Ejercicios resueltos}
\begin{ejercicio}[Lectura 18, ejercicio 1]
Si un espacio métrico tiene una cantidad finita de componentes, probar que cada componente es abierta.
\end{ejercicio}
\begin{solucion}
Sean $C_1,\ldots,C_m$ las componentes. Cada una es cerrada, y para un índice $i$,
\[
X\setminus C_i=\bigcup_{j\ne i}C_j
\]
es una unión finita de cerrados, por tanto cerrada. Su complemento $C_i$ es abierto. La finitud es esencial para este argumento: en $\Q$ hay infinitas componentes unipuntuales y ninguna es abierta.
\end{solucion}

\begin{ejercicio}[Lectura 18, ejercicio 2]
Probar que todo espacio ultramétrico es totalmente disconexo.
\end{ejercicio}
\begin{solucion}
Recordemos que una ultramétrica satisface $d(x,z)\le\max\{d(x,y),d(y,z)\}$. Toda bola abierta $B(a,r)$ es también cerrada. En efecto, si $y\notin B(a,r)$, entonces $d(a,y)\ge r$. Para $z\in B(y,d(a,y))$, si $d(a,z)<d(a,y)$, la desigualdad ultramétrica daría
\[
d(a,y)\le\max\{d(a,z),d(z,y)\}<d(a,y),
\]
una contradicción. Así $d(a,z)\ge d(a,y)\ge r$, y la bola alrededor de $y$ queda en el complemento de $B(a,r)$. El complemento es abierto.

Dados $x\ne y$, la bola $B(x,d(x,y))$ contiene a $x$ y no a $y$, y es abierta y cerrada. Sus intersecciones con cualquier conjunto que contenga ambos puntos lo separan. Por tanto ningún conexo puede contener dos puntos distintos, que es la desconexión total.
\end{solucion}

\begin{ejercicio}[Lectura 18, ejercicio 3]
Un espacio es totalmente separado si para cualesquiera puntos \emph{distintos} $x,y$ hay una separación $U,V$ del espacio con $x\in U$, $y\in V$. Probar que Cantor es totalmente separado.
\end{ejercicio}
\begin{solucion}
Mediante el homeomorfismo $C\cong\{0,1\}^{\N}$, dos puntos distintos corresponden a sucesiones que difieren en alguna coordenada $k$. Si el primero tiene dígito cero y el segundo dígito uno, usamos
\[
U=\{b:b_k=0\},\qquad V=\{b:b_k=1\}.
\]
Son abiertos en la topología producto porque la proyección de coordenada es continua y los singletons de $\{0,1\}$ son abiertos. Son complementarios y no vacíos. Transportados a $C$, dan la separación requerida. Si los dígitos están en el orden contrario, intercambiamos $U$ y $V$. La precisión ``distintos'' es necesaria: un mismo punto no puede pertenecer a dos abiertos disjuntos.\index{separacion total@separación total}
\end{solucion}

\begin{ejercicio}[Lectura 18, ejercicio 4]
Si $A\subset X$ es abierto y cerrado, probar que es unión de componentes de $X$.
\end{ejercicio}
\begin{solucion}
Si una componente $C$ toca $A$, los conjuntos $C\cap A$ y $C\cap(X\setminus A)$ son abiertos y cerrados relativos, disjuntos, y cubren $C$. La conexidad de $C$ y el hecho de que $C\cap A$ sea no vacío obligan a que el otro sea vacío. Por tanto $C\subset A$. Cada punto de $A$ pertenece a una componente, y acabamos de probar que esa componente entera está en $A$. Así $A$ es la unión de las componentes que toca. No se afirma que toda unión de componentes sea abierta: en $\Q$ un singleton es una unión de componentes pero no es abierto.
\end{solucion}

\begin{ejercicio}[Lectura 18, ejercicio 5]
Probar que $A=\{(x,x\sin(1/x)):x\ne0\}\cup\{(0,0)\}$ es arcoconexo.
\end{ejercicio}
\begin{solucion}
Definimos $h:\R\to\R$ por $h(x)=x\sin(1/x)$ para $x\ne0$ y $h(0)=0$. Fuera del cero es continua, y en cero la estimación $|h(x)|\le|x|$ prueba continuidad. El mapa $F(x)=(x,h(x))$ es continuo e inyectivo, y su inversa sobre $A$ es la primera proyección, que es continua. Por tanto $A$ es homeomorfo a $\R$. Explícitamente, para dos puntos distintos $F(a),F(b)$, el mapa
\[
\gamma(t)=F((1-t)a+tb),\qquad0\le t\le1,
\]
es continuo e inyectivo porque $a\ne b$ y la primera coordenada es estrictamente monótona. Tiene los extremos requeridos y su imagen está en $A$. El factor $x$ amortigua las oscilaciones y permite atravesar el origen, a diferencia del ejemplo \ref{ch18:seno}.
\end{solucion}

\begin{ejercicio}[Lectura 18, ejercicio 6]
Sean $X$ conexo, $A\subset X$ conexo y $B$ una componente de $X\setminus A$. Probar que $X\setminus B$ es conexo.
\end{ejercicio}
\begin{solucion}
Si $A=\varnothing$, la única componente del espacio no vacío conexo $X$ es $X$, así que $X\setminus B=\varnothing$ es conexo. Supongamos $A\ne\varnothing$ y que $Y=X\setminus B$ admite una separación $U,V$. Puesto que $A\subset Y$ es conexo, debe estar enteramente en uno de sus miembros; renombramos para que $A\subset U$. El conjunto $V$ es no vacío, abierto y cerrado en $Y=X\setminus B$.

Aplicamos el ejercicio \ref{ch17:ex6} del capítulo anterior al espacio conexo $X$, al conjunto conexo $B$ y al conjunto $V$, que es abierto y cerrado en su complemento. Obtenemos que $B\cup V$ es conexo. Como $V\cap A=\varnothing$ y $B\subset X\setminus A$, esta unión es un subconjunto conexo de $X\setminus A$. Contiene estrictamente a $B$, porque $V$ es no vacío y disjunto de $B$. Esto contradice la maximalidad de $B$ como componente de $X\setminus A$. Por tanto $Y$ no admite separación.
\end{solucion}

\begin{ejercicio}[Lectura 18, ejercicio 7: convención necesaria]
Definamos $x\approx y$ cuando $x=y$ o existe un arco de extremos $x,y$. Probar que es una relación de equivalencia y que sus clases, las componentes arcoconexas, están contenidas en componentes conexas.
\end{ejercicio}
\begin{solucion}
Sin la alternativa $x=y$, la relación escrita con arcos inyectivos no sería reflexiva: una aplicación inyectiva de $[0,1]$ no puede tener los mismos valores en $0$ y $1$. Con la convención corregida, la reflexividad forma parte de la definición. Para simetría, si $p$ es un arco de $x$ a $y$, $t\mapsto p(1-t)$ es un arco de $y$ a $x$. Para transitividad, los casos con puntos iguales son inmediatos por sustitución; si hay arcos de $x$ a $y$ y de $y$ a $z$, el lema \ref{ch18:recorte} produce un arco de $x$ a $z$ cuando son distintos, y si coinciden se aplica la primera alternativa. Así la relación es de equivalencia.

Además, si $p$ une $x$ y $y$, todo punto de su imagen está en la clase de $x$, porque la restricción de $p$ al intervalo que va de $0$ al parámetro de ese punto, reparametrizada linealmente, es un arco, salvo el punto inicial que ya está incluido por igualdad. En consecuencia, las clases son arcoconexas como subespacios, no solo pares de puntos conectables mediante arcos que podrían salir de la clase. Por la proposición \ref{ch18:arcoimplica}, cada clase es conexa; la proposición \ref{ch18:componentes} la sitúa dentro de una componente conexa. También se obtiene directamente de que la imagen de cualquier arco es un conjunto conexo que contiene sus dos extremos.\index{componente arcoconexa}
\end{solucion}

\section{Idea que conviene conservar}
Las componentes son las piezas conexas maximales y siempre son cerradas. La conexión por curvas exige un argumento adicional; la curva del seno muestra por qué. En abiertos de espacios normados, la geometría local de las bolas permite construir conexiones y recortar sus repeticiones para obtener arcos. Conexidad, conexión por caminos e inyectividad no deben confundirse en las demostraciones.
